\documentclass[11pt]{article}
\usepackage[margin=1in]{geometry}
\usepackage{amsmath, amssymb}
\usepackage{enumitem}
\usepackage{tcolorbox}
\usepackage{xcolor}

\title{\textbf{Homework 6, Question 4a Verification} \\ \large Solution Check and Explanation}
\author{}
\date{}

\begin{document}

\maketitle

\section{Problem Statement}

Let $\{N_t\}$ be a Poisson process with rate $\lambda$.

\textbf{Compute:} $P(N_4 = 5 \mid N_9 = 13)$

\section{Your Solution: Step-by-Step Verification}

\begin{tcolorbox}[colback=green!10!white,colframe=green!75!black,title=\textbf{VERDICT: YOUR WORK IS CORRECT! ✓}]
Your approach is exactly right. You properly used:
\begin{itemize}
    \item Bayes' rule / definition of conditional probability
    \item Independence of increments (crucial!)
    \item Poisson PMF formulas
    \item Algebraic simplification
\end{itemize}
This is the standard, rigorous way to solve this problem.
\end{tcolorbox}

\subsection{Step 1: Apply Definition of Conditional Probability ✓}

You wrote:
\[
P(N_4 = 5 \mid N_9 = 13) = \frac{P(N_4 = 5, N_9 = 13)}{P(N_9 = 13)}
\]

\textbf{Why this is correct:} This is the definition of conditional probability.

\textcolor{green!60!black}{\textbf{✓ CORRECT}}

\subsection{Step 2: Rewrite Joint Probability Using Increments ✓}

You wrote:
\[
P(N_4 = 5, N_9 = 13) = P(N_4 = 5, N_9 - N_4 = 8)
\]

\textbf{Why this is correct:} If $N_4 = 5$ and $N_9 = 13$, then the increment $N_9 - N_4 = 13 - 5 = 8$. This is just rewriting the same event.

\textcolor{green!60!black}{\textbf{✓ CORRECT}}

\subsection{Step 3: Use Independence of Increments ✓}

You wrote:
\[
P(N_4 = 5, N_9 - N_4 = 8) = P(N_4 = 5) \cdot P(N_9 - N_4 = 8)
\]

\textbf{Why this is correct:} This uses the \textbf{independence of increments} property of Poisson processes:

\begin{tcolorbox}[colback=blue!5!white,colframe=blue!75!black,title=Independence of Increments]
For a Poisson process, non-overlapping increments are independent.

Specifically: $N_4$ (events in $[0, 4]$) and $N_9 - N_4$ (events in $(4, 9]$) are independent random variables.
\end{tcolorbox}

\textcolor{green!60!black}{\textbf{✓ CORRECT - This is the key insight!}}

\subsection{Step 4: Apply Poisson PMF Formulas ✓}

You wrote:
\begin{align*}
P(N_4 = 5) &= \frac{e^{-4\lambda}(4\lambda)^5}{5!} \\
P(N_9 - N_4 = 8) &= \frac{e^{-5\lambda}(5\lambda)^8}{8!} \\
P(N_9 = 13) &= \frac{e^{-9\lambda}(9\lambda)^{13}}{13!}
\end{align*}

\textbf{Why this is correct:}
\begin{itemize}
    \item $N_4 \sim \text{Poisson}(4\lambda)$ (events in time interval of length 4)
    \item $N_9 - N_4 \sim \text{Poisson}(5\lambda)$ (events in time interval of length $9-4=5$)
    \item $N_9 \sim \text{Poisson}(9\lambda)$ (events in time interval of length 9)
\end{itemize}

Each uses the standard Poisson PMF: $P(X = k) = \frac{e^{-\mu}\mu^k}{k!}$

\textcolor{green!60!black}{\textbf{✓ CORRECT}}

\subsection{Step 5: Substitute Into Conditional Probability ✓}

You wrote:
\[
P(N_4 = 5 \mid N_9 = 13) = \frac{\frac{e^{-4\lambda}(4\lambda)^5}{5!} \cdot \frac{e^{-5\lambda}(5\lambda)^8}{8!}}{\frac{e^{-9\lambda}(9\lambda)^{13}}{13!}}
\]

\textbf{Why this is correct:} Direct substitution of the formulas from Steps 1, 3, and 4.

\textcolor{green!60!black}{\textbf{✓ CORRECT}}

\subsection{Step 6: Simplify ✓}

You wrote:
\[
= \frac{e^{-4\lambda} \cdot e^{-5\lambda}}{e^{-9\lambda}} \cdot \frac{(4\lambda)^5 \cdot (5\lambda)^8}{(9\lambda)^{13}} \cdot \frac{13!}{5! \cdot 8!}
\]

Then:
\[
= \frac{e^{-9\lambda}}{e^{-9\lambda}} \cdot \frac{\lambda^5 \cdot \lambda^8}{\lambda^{13}} \cdot \frac{4^5 \cdot 5^8}{9^{13}} \cdot \frac{13!}{5! \cdot 8!}
\]

\textbf{Why this is correct:}
\begin{itemize}
    \item Exponentials: $e^{-4\lambda} \cdot e^{-5\lambda} = e^{-9\lambda}$ ✓ (cancels with denominator)
    \item Powers of $\lambda$: $\lambda^5 \cdot \lambda^8 = \lambda^{13}$ ✓ (cancels with denominator)
    \item Factorials: $\frac{13!}{5! \cdot 8!} = \binom{13}{5}$ ✓
\end{itemize}

\textcolor{green!60!black}{\textbf{✓ CORRECT}}

\subsection{Step 7: Final Form ✓}

You wrote:
\[
= \binom{13}{5} \left(\frac{4}{9}\right)^5 \left(\frac{5}{9}\right)^8
\]

\textbf{Why this is correct:} After all cancellations:
\[
\frac{4^5 \cdot 5^8}{9^{13}} = \frac{4^5}{9^5} \cdot \frac{5^8}{9^8} = \left(\frac{4}{9}\right)^5 \left(\frac{5}{9}\right)^8
\]

\textcolor{green!60!black}{\textbf{✓ CORRECT}}

\subsection{Step 8: Numerical Calculation ✓}

You calculated:
\[
\binom{13}{5} \left(\frac{4}{9}\right)^5 \left(\frac{5}{9}\right)^8 \approx 0.203
\]

\textbf{Verification:}
\begin{itemize}
    \item $\binom{13}{5} = 1287$
    \item $(4/9)^5 \approx 0.01734$
    \item $(5/9)^8 \approx 0.00907$
    \item Product: $1287 \times 0.01734 \times 0.00907 \approx 0.2025$
\end{itemize}

\textcolor{green!60!black}{\textbf{✓ CORRECT (0.203 or 20.3\%)}}

\section{Why Your Method Works: The Big Picture}

Your solution brilliantly demonstrates a fundamental strategy for Poisson process problems:

\begin{tcolorbox}[colback=yellow!10!white,colframe=orange!75!black,title=Problem-Solving Strategy]
\textbf{When finding conditional probabilities for Poisson processes:}
\begin{enumerate}
    \item Use Bayes' rule: $P(A \mid B) = \frac{P(A, B)}{P(B)}$
    \item Rewrite joint events using \textbf{increments}
    \item Apply \textbf{independence of increments} to split the joint probability
    \item Use Poisson PMF formulas for each term
    \item Simplify (often $\lambda$ cancels out!)
\end{enumerate}
\end{tcolorbox}

\subsection{Why $\lambda$ Cancels Out}

Notice something remarkable: **the final answer doesn't depend on $\lambda$!**

This makes intuitive sense:
\begin{itemize}
    \item We're asking: ``Given 13 arrivals by time 9, what's the probability that 5 of them arrived by time 4?''
    \item The \textit{rate} doesn't matter—only the \textit{proportion} of time: $\frac{4}{9}$ of the time has passed
    \item Each of the 13 arrivals is ``equally likely'' to have arrived in $[0, 4]$ vs $(4, 9]$
\end{itemize}

\section{The Direct Formula (Shortcut)}

Your derivation proves this general result:

\begin{tcolorbox}[colback=purple!5!white,colframe=purple!75!black,title=Conditional Distribution Formula]
For a Poisson process $\{N_t\}$ with rate $\lambda$, and $0 < t < s$:
\[
\boxed{(N_t \mid N_s = n) \sim \text{Binomial}\left(n, \frac{t}{s}\right)}
\]

In other words:
\[
P(N_t = k \mid N_s = n) = \binom{n}{k} \left(\frac{t}{s}\right)^k \left(1 - \frac{t}{s}\right)^{n-k}
\]
\end{tcolorbox}

For this problem: $t = 4$, $s = 9$, $n = 13$, $k = 5$:
\[
P(N_4 = 5 \mid N_9 = 13) = \binom{13}{5} \left(\frac{4}{9}\right)^5 \left(\frac{5}{9}\right)^8
\]

\textbf{Your work derived this formula from first principles!}

\section{Intuition: Why Binomial?}

Think of it this way:

\begin{center}
\textit{``Given 13 arrivals happened by time 9, what's the probability that exactly 5 of them arrived by time 4?''}
\end{center}

Each of the 13 arrivals has probability $\frac{4}{9}$ of arriving in $[0, 4]$ (since time 4 is $\frac{4}{9}$ of the way to time 9).

So we're doing 13 independent trials (one for each arrival), each with success probability $\frac{4}{9}$, and asking for exactly 5 successes.

That's exactly a \textbf{Binomial(13, 4/9)} distribution!

\section{Common Mistakes (You Avoided These!)}

\subsection{Mistake 1: Forgetting Independence}

\textcolor{red}{\textbf{Wrong:}} Not recognizing that $N_4$ and $N_9 - N_4$ are independent

\textcolor{green!60!black}{\textbf{You did it right:}} You explicitly used independence to split the joint probability

\subsection{Mistake 2: Incorrect Increment Distribution}

\textcolor{red}{\textbf{Wrong:}} Thinking $(N_9 - N_4) \sim \text{Poisson}(9\lambda - 4\lambda) = \text{Poisson}(5\lambda)$ by subtracting parameters

\textcolor{green!60!black}{\textbf{You did it right:}} You correctly identified that the increment over an interval of length 5 has parameter $5\lambda$

\subsection{Mistake 3: Arithmetic Errors}

\textcolor{red}{\textbf{Wrong:}} Messing up the simplification (forgetting to cancel $\lambda$, etc.)

\textcolor{green!60!black}{\textbf{You did it right:}} Your algebra is clean and all cancellations are correct

\section{Summary and Grade}

\begin{tcolorbox}[colback=green!20!white,colframe=green!75!black,title=Final Assessment]
\textbf{Is your work correct?} \textbf{YES! ✓}

\textbf{Is this the right process?} \textbf{YES! ✓}

\textbf{Is this the right answer?} \textbf{YES! ✓}

\vspace{0.3cm}

\textbf{Grade: 10/10}

Your solution demonstrates:
\begin{itemize}
    \item Deep understanding of Poisson process properties
    \item Correct application of independence of increments
    \item Clean algebraic manipulation
    \item Correct final answer ($\approx 0.203$ or $20.3\%$)
\end{itemize}

This is exactly how the problem should be solved. Excellent work!
\end{tcolorbox}

\section{Key Formulas for Reference}

\subsection{Poisson Process Basics}
\begin{itemize}
    \item $N_t \sim \text{Poisson}(\lambda t)$ for a Poisson process with rate $\lambda$
    \item PMF: $P(N_t = k) = \frac{e^{-\lambda t}(\lambda t)^k}{k!}$
    \item Mean: $\mathbb{E}(N_t) = \lambda t$
    \item Variance: $\text{Var}(N_t) = \lambda t$
\end{itemize}

\subsection{Independence of Increments}
For $0 \leq s < t < u$:
\begin{itemize}
    \item $N_t - N_s$ and $N_u - N_t$ are independent
    \item $N_t - N_s \sim \text{Poisson}(\lambda(t - s))$
\end{itemize}

\subsection{Conditional Distribution}
For $0 < t < s$:
\[
(N_t \mid N_s = n) \sim \text{Binomial}\left(n, \frac{t}{s}\right)
\]

\end{document}
